\subsection{一个域上的一元形式幂级数环的分式域}

若 K 是一个交换域，我们将用 \(\mathrm{K}((\mathbf{X}))\) 表示整环 \(\mathrm{K}[[\mathbf{X}]]\) 的分式域。

命题 12. --- K((X)) 的每个非零元素 u 都可以唯一地写成 \(u = X^{k}v\) ，其中 \(k \in Z\) 且 v 是 X 的 0 阶形式幂级数。设 u = w/t，其中 w, t 是 K{[}{[}X{]}{]} 的非零元素。我们有 \(w = X^{r}w_{1}, t = X^{s}t_{1}\) ，其中 r, \(s \in N\) 且 \(w_{1}, t_{1}\) 是 0 阶形式幂级数，从而在 K{[}{[}X{]}{]} 中可逆（IV，第30页，命题 6）。于是 \(u = X^{r-s}w_{1}t_{1}^{-1}\) ，且 \(w_{1}t_{1}^{-1}\) 是 0 阶形式幂级数。

我们来证明唯一性。假设 \(u = X^{k_{1}}v_{1} = X^{k_{2}}v_{2}\) ，其中 \(k_{1}, k_{2} \in Z\) 且 \(v_{1}, v_{2}\) 是 0 阶形式幂级数。由于 \(X^{k_{1}-k_{2}} = v_{2}v_{1}^{-1}\) 是 0 阶形式幂级数，我们有 \(k_{1} = k_{2}\) ，由此 \(v_{1} = v_{2}\) ，这就证明了唯一性断言。

我们将称 \(\mathbf{K}((\mathbf{X}))\) 中的元素为 X 中系数在 K 内的广义形式幂级数，或者在没有混淆时简称为形式幂级数（ \(K[[X]]\) 中的元素则称为具有正指数的形式幂级数）；若 \(u \neq 0\) ，命题 12 中定义的整数 k 也称为 u 的阶，并记作 \(\omega(u)\) ，即使它是 \textless{} 0 时也如此；我们还记 \(\omega(0) = \infty\) 。可以立即验证

\[
\begin{array}{l} \omega (u + v) \geqslant \inf (\omega (u), \omega (v)) \\ \omega (u + v) = \inf (\omega (u), \omega (v)) \quad \text { if } \quad \omega (u) \neq \omega (v) \\ \omega (u v) = \omega (u) + \omega (v) \end{array}
\]

对广义形式幂级数仍然成立。特别地，若 \(u \neq 0\) ，则 \(\omega(u^{-1}) = -\omega(u)\) 。* 换言之（《交换代数》，VI，第 3 节，第 6 号，第 392 页，定义 3），\(\omega\) 是域 \(K((X))\) 的一个规范化离散赋值。*

对每个整数 \(n \in \mathbf{Z}\) ，令 \(\mathfrak{p}_n\) 为所有 \(u \in K((X))\) （其中 \(\omega(u) \geqslant n\) ）组成的集合。则 \((\mathfrak{p}_n)_{n \in \mathbf{Z}}\) 是加法群 \(K((X))\) 的一个递减子群序列，其交为 0；因此在 \(K((X))\) 上存在一个在平移下不变的拓扑，使得 \((\mathfrak{p}_n)_{n \in \mathbf{Z}}\) 是 0 的基本邻域系（《一般拓扑学》，III，第223页）。我们容易验证 \(K((X))\) 是一个拓扑域（《一般拓扑学》，III，第281页），并且 \(K[[X]]\) 是 \(K((X))\) 的开且闭子空间。

设 \((\alpha_{n})_{n\in \mathbf{Z}}\) 是 \(\mathbf{K}\) 中元素的一个族，并假设存在整数 \(\mathbf{N}\) 使得 \(\alpha_{n} = 0\) 对所有 \(n < \mathbf{N}\) 。则族 \((\alpha_{n}\mathbf{X}^{n})_{n\in \mathbf{Z}}\) 在 \(\mathbf{K}((\mathbf{X}))\) 中可和（《一般拓扑学》，III，第263页，推论）；令 \(u = \sum_{n\in \mathbf{Z}}\alpha_n\mathbf{X}^n\) ，则 \(u = 0\) 当且

仅当 \(\alpha_{n} = 0\) 对所有 \(n\) ；否则 \(u\) 的阶是最小整数 \(k\) ，它使 \(\alpha_{k} \neq 0\) 。最后，\(\mathbf{K}((\mathbf{X}))\) 的每个元素都可以唯一地写成 \(\sum_{n \in \mathbf{Z}} \alpha_{n} \mathbf{X}^{n}\) 的形式，其中序列 \((\alpha_{n})\) 满足 \(\alpha_{-n} = 0\) 对所有充分大的 \(n\) 。

由于环 K{[}X{]} 是 K{[}{[}X{]}{]} 的子环，每个有理分式 \(u/v \in \mathrm{K}(\mathrm{X})\) （u, v 为 X 的多项式）都可以等同于 K((X)) 中的（广义）形式幂级数 \(uv^{-1}\) ，我们称之为它在原点的展开式；域 K(X) 由此被等同于 K((X)) 的一个子域。

